% ============================================================
% CONTROL DEL AVANCE Y COMANDOS PROPIOS
% ============================================================

% Seleccionar un unico valor entero:
% 1 = Avance 1
% 2 = Avance 2 / primera entrega
% 3 = Avance 3
% 4 = Avance 4 / paper final
\providecommand{\avanceActual}{4}

% Imprime el segundo argumento desde el avance indicado en adelante.
% Ejemplo: \incluirDesdeAvance{3}{contenido visible desde el avance 3}
\newcommand{\incluirDesdeAvance}[2]{%
  \ifnum\avanceActual<#1\relax
  \else
    #2%
  \fi
}

% ============================================================
% AJUSTES APA 7 COMPATIBLES CON ELSARTICLE
% ============================================================

% Aplicar después del frontmatter para no alterar portada, título, autores,
% afiliaciones, resumen ni palabras clave.
\newcommand{\activarFormatoCuerpoAPA}{%
  \setlength{\parindent}{1.27cm}%
  \setlength{\parskip}{0pt}%
}

% biblatex usa \bibhang para la sangría francesa de cada referencia.
\newcommand{\prepararReferenciasAPA}{%
  \setlength{\parindent}{0pt}%
  \setlength{\bibhang}{1.27cm}%
  \setlength{\bibitemsep}{0pt}%
}

% Infraestructura recomendada para tablas:
% \begin{threeparttable}
%   \caption{Título claro de la tabla}
%   \begin{tabular}{...} ... \toprule ... \bottomrule \end{tabular}
%   \begin{tablenotes}\item \textit{Nota.} ... \end{tablenotes}
% \end{threeparttable}
% No usar líneas verticales ni sombreado salvo justificación académica.
